\documentclass[12pt]{article}
\usepackage[utf8]{inputenc}
\usepackage[T1]{fontenc}
\usepackage{hyperref}
\title{Moi Lemas Language Resources for K-12 Curriculum Development}
\author{Compiled by Claude Sonnet 5.5 (claude-sonnet-5-5) for Andr\'as Kornai}
\date{}
\begin{document}
\maketitle

\section{Introduction}
Moi Lemas (also Seget, Moi Sigin; ISO 639-3 sbg, Glottocode sege1235) is a West Bird's Head language of Sorong Regency, Southwest Papua, Indonesia. Villages named in sources include Walian, Sailolof, Segum and Seget. About 1,200 speakers are cited, from a 1988 figure. It should not be confused with Moi (mxn), a separate language. Searches found no school materials and no published dictionary, grammar or Scripture. Only small documentation collections and general reference pages turned up.

\section{Primary Educational Resources}
No primers, textbooks or school programmes in Moi Lemas were found in search results.

\section{Linguistic and Reference Materials}
\begin{itemize}
\item Seget language (Wikipedia overview): \url{https://en.wikipedia.org/wiki/Seget_language}
\item West Bird's Head languages (Wikipedia): \url{https://en.wikipedia.org/wiki/West_Bird%27s_Head_languages}
\item University of Edinburgh collection by Laura Arnold (2023): elicited morphosyntactic data from two speakers, including possessive paradigms, with audio and field notes. \url{https://datashare.ed.ac.uk/handle/10283/8695} and \url{https://www.research.ed.ac.uk/en/datasets/seget-iso-369-3-sbg-glottocode-sege1235-collection/}
\end{itemize}

\section{Religious and Community Resources}
\begin{itemize}
\item Joshua Project people-group profile (Seget): \url{https://joshuaproject.net/people_groups/14817}
\end{itemize}
No Scripture translation was found.

\section{Digital and Audio Resources}
\begin{itemize}
\item Global Recordings Network lists varieties ``Moi Lemas: Yauw'' and ``Moi Lemas: Moi Sigin'': \url{https://globalrecordings.net/en/language/29176}
\end{itemize}

\section{Curriculum Development Context}
The language is spoken by adults and is not known to be taught in schools. Very little documentation exists, so any curriculum would depend on new community-led documentation. The Edinburgh data could serve as a starting point for a phonological and morphological sketch.
\end{document}